\chapter{Diseño de la investigación y métodos}
\label{ch:methods}
\section{Diseño y justificación}
Explica por qué el diseño responde a la pregunta de investigación. Identifica
la unidad de análisis, las condiciones de comparación y los controles.
Describe la selección de participantes y los criterios de inclusión y exclusión.
\section{Materiales y procedimiento}
Documenta los materiales, las herramientas y sus versiones. Describe cada fase
con suficiente detalle para reproducirla.
\begin{table}[htbp]
\centering
\caption{Protocolo de ejemplo; adáptalo a tu investigación.}
\begin{tabularx}{\linewidth}{@{}lX@{}}
\toprule
\textbf{Fase} & \textbf{Información que registrar}\\
\midrule
Preparación & Tarea, materiales y criterios de selección.\\
Recogida & Condiciones, observaciones y desviaciones del protocolo.\\
Análisis & Métricas, incertidumbre, exclusiones y código de análisis.\\
\bottomrule
\end{tabularx}
\end{table}
\section{Medidas y plan de análisis}
Define los resultados, las unidades y la comparación antes de presentar los datos.
Para el tiempo de realización, una posible medida es
\begin{equation}
\Delta_t = \frac{t_{\mathrm{habitual}}-t_{\mathrm{revisado}}}{t_{\mathrm{habitual}}}\times100\%.
\end{equation}
Justifica el análisis estadístico o cualitativo y el tratamiento de datos ausentes.
Distingue los análisis confirmatorios de los exploratorios.
\section{Ética y reproducibilidad}
Describe el consentimiento, la privacidad, la gestión de datos y las autorizaciones.
Indica cómo compartirás los materiales, datos y código, o justifica las restricciones.
